% =====================================================================
%  CHAPTER 5 - IMPLEMENTATION
% =====================================================================
\clearpage
\chapter{Implementation}

The engine was implemented in Go~1.25 as a single binary, \cd{engined}, of
about 43,000 lines of non-test code, together with roughly 23,000 lines of
tests. This chapter describes the structure of the code and the parts of the
implementation where the details matter for correctness or performance. The
design that the code realises was given in Chapter~4 and is not repeated.

\section{Code Organisation}

Table~\ref{tab:packages} lists the main packages. Dependencies point
downwards: \cd{core} defines the event types, identifiers and interfaces and
imports nothing from the engine, and the storage packages depend only on
\cd{core} and the PostgreSQL driver. The runtime package wires the pieces
together, and \cd{cmd/engined} builds the process from configuration.

\begin{table}[htbp]
  \centering
  \caption{Main packages of the engine}
  \label{tab:packages}
  \small
  \begin{tabular}{@{}p{1.25in}p{0.55in}p{3.7in}@{}}
    \toprule
    \textbf{Package} & \textbf{Lines} & \textbf{Responsibility} \\
    \midrule
    \cd{core} & 625 & Event types, identifiers, decision and store interfaces \\
    \cd{eventlog} & 1,005 & Reading histories; the fenced decision commit \\
    \cd{taskq} & 1,002 & Durable task queue: poll, lock, heartbeat, complete,
      retry \\
    \cd{lease} & 789 & Lease store and the ownership state machine \\
    \cd{shard} & 275 & Workflow-to-shard hash and the consistent-hash ring \\
    \cd{timers} & 264 & Durable timers and the due-timer scanner \\
    \cd{workflow} & 1,644 & Workflow context, replay by program order,
      cooperative dispatcher \\
    \cd{activity} & 729 & Activity registry, retry policy, deduplication keys \\
    \cd{saga} & 202 & Compensation stack \\
    \cd{idem} & 3,454 & Exact and semantic deduplication, embedder client,
      hybrid comparator \\
    \cd{runtime} & 2,954 & Workflow, activity and timer workers; per-shard
      worker pools \\
    \cd{api} & 1,959 & HTTP API, error mapping, request routing \\
    \cd{obs} & 444 & Prometheus metrics and structured logging \\
    \cd{faultinj} & 364 & Named crash points for tests \\
    \bottomrule
  \end{tabular}
\end{table}

\section{The Fenced Commit}

All writes produced by one workflow task are made by a single function,
\cd{CommitDecision}, in one transaction. Its first statement is

\begin{lstlisting}[caption={Fencing check at the start of every commit},
  label={lst:verify}]
SELECT fencing_token FROM shard_leases WHERE shard_id = $1 FOR UPDATE
\end{lstlisting}

\noindent and the function returns a \cd{NotOwner} error, after rolling back,
if the value differs from the token held by the caller. The remaining steps
follow the order given in Section~4.2.2. Placing the start of a new run
after the fencing check was a deliberate change made during development: in
an earlier version the insert into \cd{executions} for a start was performed
outside the fenced transaction, so a node that had lost a shard could still
admit new runs to it. The same pattern, a token check under a row lock in the
writing transaction, is used by the timer scanner and by the signal path.

\section{Task Queue}

Each worker polls one owned shard at a time with a query that selects one
ready row and locks it in the same statement:

\begin{lstlisting}[caption={Task poll (simplified)}, label={lst:poll}]
WITH pick AS (
  SELECT shard_id, task_id FROM tasks
   WHERE shard_id = $1 AND task_type = $2 AND ready_at <= now()
   ORDER BY shard_id, task_type, ready_at
   FOR UPDATE SKIP LOCKED LIMIT 1)
UPDATE tasks t
   SET locked_until = now() + $3::interval,
       lock_token = gen_random_uuid(), attempt = t.attempt + 1
  FROM pick p
 WHERE t.shard_id = p.shard_id AND t.task_id = p.task_id
RETURNING t.*
\end{lstlisting}

\noindent Here \cd{ready\_at} stands for the later of the visibility time and
the lock expiry.
A partial index on ready tasks, added in the second migration, keeps the
inner query to an index scan. Completion and failure are accepted only with
the lock token issued by the poll, so a worker whose lock has expired and
whose task has been given to another worker cannot complete it. Idle workers
back off between polls; Chapter~7 shows that this polling is the dominant
database load on an idle cluster and that it sets part of the latency floor.

\section{Lease Manager}

The ownership policy is written as two pure functions. \cd{Decide} takes the
remembered state of one shard and the current time and returns one action:
acquire, renew, declare lost or do nothing. \cd{Apply} takes the outcome of
the database call and returns the new state. Neither function touches the
database or a clock, so all combinations of late renewals, failed queries and
expired leases are covered by table-driven unit tests that run in
microseconds. Two rules in these functions matter for safety. Expiry is
checked before renewal, so a node never renews a lease that has already
lapsed. A renewal that fails because the database is unreachable does not
advance the time of the last successful renewal, so if the failure persists
the node gives up the shard on its own. The loop that calls these functions
runs every renewal interval and calls back into the runtime to start or stop
the shard's workers.

\section{Workflow Execution}

Each workflow task runs the workflow function in a goroutine controlled by a
dispatcher. When the function calls \cd{Get} on a future that has no recorded
result, the goroutine hands control back and is then discarded; the run will
be re-executed from the start on its next task. This keeps the runtime free of
long-lived per-run goroutines, so a run sleeping for a week costs one row in
\cd{timers} and nothing in memory. The context counts calls in program order
and, on each call, compares the requested command with the event recorded at
that position; a mismatch fails the run with a non-determinism error that
names the position, the recorded command and the requested one.

\section{Semantic and Hybrid Deduplication}

The semantic checker builds a canonical intent string from the activity type
and its normalised, sorted arguments, obtains an embedding from the embedder
service, and compares it by cosine similarity with the stored embeddings of
completed activities of the same type in the run. The verdict and its inputs
are journalled as a \cd{SemanticDedupEvaluated} event before any command is
issued. Lookups and inserts of embeddings use the \cd{pgvector} type in
\cd{activity\_intents}.

The evaluation in Chapter~7 showed that a general sentence encoder answers a
different question from the one deduplication asks. It judges two requests
for the same route on different days as very similar, which they are, though
they are two different bookings. An opt-in \emph{hybrid} comparator was
therefore added. The activity declares its deciding fields, for example the
origin, destination and date of a flight; these are normalised (``BLR'' and
``Bangalore'' map to one code, ``Friday AM'' to a date and time window) and
compared exactly, and similarity is applied only to the remaining free text.
The verdict is journalled in the same way, so the hybrid comparator changes
the decision but not the replay mechanism.

\section{API, Routing and Observability}

The HTTP API (Table~\ref{tab:api}) is served by every node. A request for a
run on a shard the node does not own is answered with HTTP~409 \cd{NotOwner} and, when
known, the name of the current owner, and clients retry against that node. Every node exports
Prometheus metrics for task latencies, commits, lease transitions, fencing
rejections, deduplication verdicts and connection-pool use, and writes
structured JSON logs. The fencing-rejection counter is the one used in
Chapter~7 to count stale writes.

\begin{table}[htbp]
  \centering
  \caption{HTTP API}
  \label{tab:api}
  \small
  \begin{tabular}{@{}p{2.6in}p{2.9in}@{}}
    \toprule
    \textbf{Endpoint} & \textbf{Purpose} \\
    \midrule
    \cd{POST /v1/workflows} & Start a run (idempotent on identifier) \\
    \cd{GET /v1/workflows/\{id\}} & Status of a run \\
    \cd{GET /v1/workflows/\{id\}/history} & Event history of a run \\
    \cd{POST /v1/workflows/\{id\}/signals} & Send a named signal \\
    \cd{POST /v1/workflows/\{id\}/cancel} & Request cancellation \\
    \cd{GET /v1/shards} & Current lease table \\
    \cd{GET /v1/dedup/verdicts} & Recorded semantic verdicts \\
    \cd{GET /healthz}, \cd{GET /metrics} & Liveness and metrics \\
    \bottomrule
  \end{tabular}
\end{table}

\section{Demonstration Applications}

Two applications exercise the engine. The first is a food-delivery order
written as a saga: payment, restaurant confirmation, driver assignment and
notification, each with a compensation, and a durable timer for the
preparation time. Every step has a natural exact key such as the order
identifier, so it uses no semantic deduplication. It is the workload used for
all load tests in Chapter~7. The second is a travel agent whose planner
proposes flight and hotel bookings. It can run with a recorded fixture or a
live language model, and it reproduces the double-booking case of
Section~1.2.3 when a crash is injected between two planning steps.

\section{Deployment}

The binary is built into a small container image. Locally, Docker Compose
starts three engine nodes, PostgreSQL with \cd{pgvector}, the embedder,
Prometheus and Grafana. On Kubernetes the engine runs as a StatefulSet, so
that each pod has a stable name that is also its node identifier on the ring;
pod anti-affinity places one engine per node. The AWS environment of
Chapter~6 is created by Terraform and contains no component other than
those listed here: no message broker, cache or coordination service.
